% ----------------   DA-Inhalt Deutsch
\newcommand*{\topicG} 	{Setzen der Diplomarbait in \LaTeX}
\newcommand*{\partner}	{TeX Users Group}
\newcommand*{\purposeG}%
	{Erstellen eines \LaTeX -Vorlagendokuments für Diplomarbeiten an der HTL Weiz.
	 Das Dokument soll sowohl die Sturktur als auch das Arbeiten mit \LaTeX ~beschreiben.}
\newcommand*{\implementationG}%
	{Das Dokument soll technisch versierten jedoch in \LaTeX ~ungeübten StudentInnen dabei helfen, ihre Diplomarbeit mit \LaTeX ~zu erstellen.
	 Dabei soll sowohl eine ausreichend konfigurierte Vorlage, als auch Informationen zur Verwendung dieser bereit gestellt werden.
	 Auf eine korrekte Rechtschreibung wurde weniger geachetet.}
\newcommand*{\resultsG}%
	{Als Ergbenis steht diese Vorlage zur Verwendung zur Verfügung.
	 Sie darf gerne abgeändert und\slash oder erweitert werden.
	 Rückmeldung an \href{mailto:qui@htlweiz.at}{qui@htlweiz.at} werden begrüßt!}
\newcommand*{\illustrativeG}{../00_DefTex/HtlHintergrundA4_2.png}		%Pfadangabe mit Dateinamen zum Bild der Kurzfassung
\newcommand*{\participationsG}	{HTL-Weiz \LaTeX -Convention}
\newcommand*{\accessibilityG} 	{Zur freien Einsicht freigegeben.}

% ----------------   DA-Inhalt Englisch
\newcommand*{\topicE} 	{Typesetting the diploma thesis in \LaTeX}
\newcommand*{\purposeE}%
	{Creating a \LaTeX template document for diploma theses at the HTL Weiz.
	 The document should describe both the structure and how use \LaTeX .}
\newcommand*{\implementationE}%
	{The document is intended to help technically experienced students who are not familiar with \LaTeX to create their diploma thesis in \LaTeX.
	 Both, a sufficiently configured template and information on how to use it should be provided.
	 Less attention was paid to correct spelling.}
\newcommand*{\resultsE}%
	{As a result, this template is available for free use.
	 You are welcome to be modifie and extend.
	 Feedback to \href{mailto:qui@htlweiz.at}{qui@htlweiz.at} is appreciated!}
\newcommand*{\illustrativeE}{\illustrativeG}							%Kan so gelassen werden Wenn das Bild in Deutsch und Englisch das selbe ist
\newcommand*{\participationsE}	{HTL-Weiz \LaTeX - Convention}
\newcommand*{\accessibilityE} 	{Released for free viewing}